\documentclass[12pt,a4paper]{article}
\usepackage[utf8]{vietnam}
\usepackage{amsmath,amssymb}
\usepackage{tabularx}
\usepackage{tikz}
\usepackage{xcolor}
\usepackage[top=0.6cm, bottom=0.6cm, left=0.6cm, right=0.6cm]{geometry}
\usepackage[dethi]{ex_test}

\SAOPTN{widthSA=2.8,heightSA=0.62}
\setlength{\parskip}{0pt}

% Dòng kẻ chấm mờ chuẩn viết tay
\newcommand{\dongke}[1]{%
  \begingroup%
  \setlength{\parskip}{0pt}%
  \baselineskip=18.5pt%
  \vspace{1.0mm}%
  \foreach \i in {1,...,#1}{%
    \noindent{\color{gray!60}\dotfill}\par%
  }%
  \vspace{0.6mm}%
  \endgroup%
}

% Định nghĩa ô bài tập tự luận trong bảng
\newcommand{\oBaiTap}[3][6]{%
  \vspace{1.0mm}%
  \noindent\textbf{#2} #3\par%
  \dongke{#1}%
}

\begin{document}
\pagestyle{empty}
% =============================================================================
% ========================= MA DE 05 - MAT TRUOC (TRANG 1/2) =============
% =============================================================================
\noindent
\begin{tabularx}{\linewidth}{|X|c|}
\hline
\textbf{PHIẾU KIỂM TRA ĐỊNH KỲ TOÁN 6 - TUẦN 12 - BUỔI 1} & \textbf{Mã đề: 05 - Trang 1/2} \\
\hline
\end{tabularx}

\vspace{1.0mm}

\noindent
\begin{tabularx}{\linewidth}{|p{3.2cm}|X|}
\hline
\centering\textbf{Điểm số} & \centering\textbf{Lời nhận xét của thầy cô} \tabularnewline
\rule{0pt}{1.8cm} & \\
\hline
\end{tabularx}

\vspace{1.0mm}
\noindent\textbf{Câu 1 (0,25 điểm).} Bỏ dấu ngoặc của biểu thức $-(x - 15)$, ta được kết quả là: \shortans{-x + 15}\par\vspace{0.10mm}
\noindent\textbf{Câu 2 (0,25 điểm).} Bỏ dấu ngoặc của biểu thức $+(-28 + x)$, ta được kết quả là: \shortans{-28 + x}\par\vspace{0.10mm}
\noindent\textbf{Câu 3 (0,25 điểm).} Giá trị của biểu thức $(-58) + 32 + 58$ là: \shortans{32}\par\vspace{0.10mm}
\noindent\textbf{Câu 4 (0,25 điểm).} Bỏ dấu ngoặc của biểu thức $-(u - v + 6)$, ta được kết quả là: \shortans{-u + v - 6}\par\vspace{0.10mm}
\noindent\textbf{Câu 5 (0,25 điểm).} Điền vào chỗ trống: $50 - u - v = 50 - (\dots)$ là: \shortans{u + v}\par\vspace{0.10mm}
\noindent\textbf{Câu 6 (0,25 điểm).} Điền dấu thích hợp ($<, >, =$) vào chỗ trống: $-(-32) + (-40) \; \dots \; 0$ \shortans{<}\par\vspace{0.10mm}
\noindent\textbf{Câu 7 (0,25 điểm).} Hình thoi có số trục đối xứng là: \shortans{2}\par\vspace{0.10mm}
\noindent\textbf{Câu 8 (0,25 điểm).} Trong các chữ cái X, V, T, Y, chữ cái có tâm đối xứng là: \shortans{X}\par\vspace{0.10mm}
\noindent\textbf{Câu 9 (0,25 điểm).} Hình chữ nhật có chiều dài $14\text{ cm}$, chiều rộng $5\text{ cm}$. Diện tích là: \shortans{70\text{ cm}^2}\par\vspace{0.10mm}
\noindent\textbf{Câu 10 (0,25 điểm).} Nhiệt độ trưa $26^\circ\text{C}$, đêm giảm xuống $17^\circ\text{C}$. Độ chênh lệch nhiệt độ là: \shortans{9^\circ\text{C}}\par\vspace{0.10mm}

\vspace{0.3mm}
\noindent\textbf{Câu 11 (3,0 điểm). Thực hiện phép tính (hợp lý nếu có thể):}
\vspace{0.4mm}

\noindent
\begin{tabularx}{\linewidth}{|X|X|}
\hline
\oBaiTap[7]{a) (0,5 điểm)}{$165 - (65 - 165)$} & \oBaiTap[7]{b) (0,5 điểm)}{$(-95) - (25 - 190)$} \\
\hline
\oBaiTap[7]{c) (0,5 điểm)}{$(482 - 76) - (482 - 176)$} & \oBaiTap[7]{d) (0,5 điểm)}{$190 - [75 - (20 - 55)]$} \\
\hline
\oBaiTap[7]{e) (0,5 điểm)}{$60 - 4 \cdot (-2)^3$} & \oBaiTap[7]{f) (0,5 điểm)}{$(27 - 38 + 65) - (27 - 38 - 35)$} \\
\hline
\end{tabularx}
\newpage
% =============================================================================
% ========================= MA DE 05 - MAT SAU (TRANG 2/2) ==============
% =============================================================================
\noindent
\begin{tabularx}{\linewidth}{|X|}
\hline
\centering\textbf{TOÁN 6 - TUẦN 12 - BUỔI 1 | Mã đề: 05 - Trang 2/2} \\
\hline
\end{tabularx}

\vspace{1.2mm}

\noindent\textbf{Câu 12 (2,0 điểm). Tìm số nguyên $x$, biết:}
\vspace{0.5mm}

\noindent
\begin{tabularx}{\linewidth}{|X|X|}
\hline
\oBaiTap[5]{a) (0,5 điểm)}{$x - 16 = -35$} & \oBaiTap[5]{b) (0,5 điểm)}{$36 - (x + 14) = 58$} \\
\hline
\oBaiTap[6]{c) (0,5 điểm)}{$28 - (18 - x) = -16$} & \oBaiTap[6]{d) (0,5 điểm)}{Liệt kê và tính tổng các số nguyên $x$ thỏa mãn: $-4 < x \le 4$} \\
\hline
\end{tabularx}

\vspace{1.2mm}

\noindent
\begin{tabularx}{\linewidth}{|X|}
\hline
\vspace{0.4mm}
\textbf{Câu 13 (2,0 điểm).} Một công ty may mặc xuất khẩu theo dõi kết quả kinh doanh trong năm và ghi nhận lợi nhuận qua từng quý (số tiền lãi ghi số dương, số tiền lỗ ghi số âm) như sau:\par\vspace{0.4mm}\centerline{\begin{tabular}{|l|c|c|c|c|}\hline \textbf{Quý trong năm} & Quý 1 & Quý 2 & Quý 3 & Quý 4 \\\hline \textbf{Lợi nhuận} & $-50\text{ triệu đồng}$ & $+140\text{ triệu đồng}$ & $-70\text{ triệu đồng}$ & $+170\text{ triệu đồng}$ \\\hline\end{tabular}}\vspace{0.4mm}\par\vspace{0.3mm}
\textbf{a) (1,0 điểm)} Tính tổng lợi nhuận của công ty trong nửa đầu năm (Quý 1 và Quý 2).\par\vspace{0.3mm}
\textbf{b) (1,0 điểm)} Tính tổng lợi nhuận của cả năm và cho biết cả năm công ty lãi hay lỗ bao nhiêu tiền.\par
\vspace{0.4mm}
\noindent\textit{Bài giải:}\dongke{9}
\\ \hline
\end{tabularx}

\vspace{1.2mm}

\noindent
\begin{tabularx}{\linewidth}{|X|}
\hline
\vspace{0.4mm}
\textbf{Câu 14 (0,5 điểm).} Tìm tất cả các số nguyên $x$ để phân số sau nhận giá trị là một số nguyên: $A = \dfrac{6}{x - 3}$.\par
\vspace{0.4mm}
\noindent\textit{Bài giải:}\dongke{8}
\\ \hline
\end{tabularx}
\end{document}
